% =====================================================================
%  26.1  Construction -- close a marked set under LCA and link it with a
%  monotone stack.  Exports: sec:c30-construction  def:c30-vt
%  thm:c30-size  fig:c30-build  lst:c30-build  prob:c30-ancient
%  lst:c30-ancient
%  Solved CF 2127E "Ancient Tree" (2100, L5), live-verified against the CF
%  problem page + the Round 1041 editorial (which explicitly builds, per colour,
%  A = S_c u {lca(u,v)} and counts how many colour virtual trees each vertex is a
%  branch of).  Official 4-test sample read live; test 1 (4-node star) replayed
%  here to cost 0.  DISCHARGES the Ch.16 (sec:c11-lifting) forward promise "a virtual tree
%  compresses a k-node set via sort-by-tin + LCA of adjacent pairs"; LCA and tin
%  are USED, never re-derived.
% =====================================================================
\section{Building a Virtual Tree}\label{sec:c30-construction}

\begin{previouslybox}[Ch.~16 handed you \(\lca\) and promised this chapter]
Chapter~16 (\(\secref{c11-lifting}\)) proved that one DFS fixes \(\tin,\tout,\depth\), that
\(u\) is an ancestor of \(v\) iff \(\tin(u)\le\tin(v)<\tout(u)\), and that binary
lifting answers \(\lca(u,v)\) -- hence \(\dist(u,v)=\depth(u)+\depth(v)-2\depth(\lca)\)
(\(\thmref{c11-distf}\)) -- in \(O(\log n)\).  It then promised, in its own words,
that ``a virtual tree (Ch.~26) compresses a \(k\)-node set via sort-by-tin plus the
LCA of adjacent pairs; this chapter's distance formula and lifting are the inner
loop.''  This section is that compression.  We never rebuild a lifting table and
never re-prove the distance formula; we call \(\lca\) \(O(k)\) times and spend the
result on a set of \(k\ll n\) marked vertices.
\end{previouslybox}

\begin{topicintro}
Given the fixed tree \(T\) and a marked set \(S\) of size \(k\), the virtual tree is
the smallest subtree-shaped \emph{view} of \(T\) that (i) contains every marked node,
(ii) is closed under \(\lca\) (so all the branching is preserved), and (iii) replaces
every maximal chain of unmarked, non-branching nodes by one weighted edge.  Building
it has exactly two steps.  \emph{Close}: sort \(S\) by \(\tin\), add \(\lca\) of each
adjacent pair, dedup.  \emph{Link}: sweep the closed set by \(\tin\) with a stack of
the current root-to-node path, hanging each node on its nearest kept ancestor.  The
miracle, proved in \thmref{c30-size}, is that \(2k-1\) nodes is a hard ceiling -- the
structure is linear in the marked set, not in \(n\).

\textbf{Prerequisites.} Chapter~16 (\(\secref{c11-paths}\)) -- \(\lca\), \(\dist\).
Chapter~5 (\(\secref{c24-times}\), \chref{c24}) -- the \(\tin/\tout\) order and the
ancestor test, used as the sort key and the stack predicate.  Nothing else.

\textbf{After this section you can}
\begin{itemize}[nosep]
  \item close a marked set under \(\lca\) with \(k-1\) adjacent-pair calls; \hfill\lv{3}
  \item prove the closed set has at most \(2k-1\) nodes; \hfill\lv{4}
  \item build the compressed tree with one monotone-stack sweep in \(\bigO{k\log k}\); \hfill\lv{4}
  \item recognise a per-colour / per-subset problem and collapse it -- solved CF 2127E. \hfill\lv{4}
\end{itemize}
\end{topicintro}

% ---------------------------------------------------------------
\subsection{Closing the marked set under \(\lca\)}
% ---------------------------------------------------------------

\begin{definition}{Virtual (auxiliary) tree of a marked set}{c30-vt}
Let \(S\subseteq V(T)\) be the marked set.  The \term{virtual tree} \(A(S)\) is the
set \(A(S)=S\cup\{\lca(x,y):x,y\in S\}\) made into a tree by drawing an edge
\((u,v)\) whenever \(u\) is a proper ancestor of \(v\) in \(T\) and no element of
\(A(S)\) lies strictly between them on the real path.  Equivalently, \(A(S)\) is the
minimal \(\lca\)-closed superset of \(S\), and it is exactly the union of all simple
paths of \(T\) between pairs of marked nodes, with maximal unbranched chains
compressed into single weighted edges.
\end{definition}

Naively \(|A(S)|\) looks like it could be as large as the whole tree, because there
are \(\binom{k}{2}\) pairs.  The load-bearing fact is that only \emph{adjacent}
pairs in \(\tin\)-order are needed to close the set.

\begin{theorem}{Adjacent-pair LCAs suffice, and \(|A(S)|\le 2k-1\)}{c30-size}
Sort \(S\) by entry time as \(x_1,x_2,\dots,x_k\).  Then
\[
   S\cup\{\lca(x_i,x_{i+1}): 1\le i<k\}
\]
is already closed under \(\lca\) -- it equals \(A(S)\).  In particular \(|A(S)|\le
2k-1\).
\end{theorem}
\begin{proofbox}
Take any two marked nodes \(a,b\) and let \(w=\lca(a,b)\).  In the DFS order the
marked nodes in \(w\)'s subtree form a contiguous block, and \(a,b\) being under
\(w\) in different child-subtrees of \(w\) (or one being \(w\)) means there is some
\emph{adjacent} pair \(x_i,x_{i+1}\) inside that block whose paths leave \(w\)
through two different children (or one is \(w\) itself); for that pair
\(\lca(x_i,x_{i+1})=w\).  Hence every \(\lca\) of marked nodes appears as an
adjacent-pair \(\lca\), so the \(k-1\) adjacent LCAs already contain all of
\(A(S)\setminus S\).  For the size: the compressed tree has \(k\) marked leaves
(every element of \(S\) that is not an LCA of others is a leaf; add back those that
are internal, at most \(k\)) and an \(\lca\)-closed tree in which every internal node
has at least two children drawn from the kept set, so the number of internal nodes
is at most \(k-1\); together \(|A(S)|\le k+(k-1)=2k-1\). \(\Box\)
\end{proofbox}

\begin{insightbox}[\(\tin\)-adjacency is doing the work of \(\binom{k}{2}\) pairs]
A branching node of the marked set is, by definition, an LCA where the marked nodes
below it split into two different DFS child-buckets.  Sorting by \(\tin\) lines those
buckets up consecutively, so the split is always visible between two \emph{neighbours}
in the order -- which is why \(k-1\) LCA calls, not \(\binom{k}{2}\), close the set.
The bound \(2k-1\) is tight: take \(k\) leaves all sharing one root; their LCAs collapse
to the single root, giving \(k+1\); take a caterpillar of nested pairs and you approach
\(2k-1\).
\end{insightbox}

% ---------------------------------------------------------------
\subsection{Linking with a monotone stack}
% ---------------------------------------------------------------

Once the closed set \(C=A(S)\) is known and sorted by \(\tin\), each node must be
hung on its nearest kept ancestor.  Walking up with \(\lca\) per node is \(O(k\log
n)\) and fiddly; the clean way is one stack sweep that maintains the current
root-to-node path inside \(C\).

\begin{lstlisting}[caption={Build the virtual tree of a marked set. Input: marked nodes, and lca/tin from a preprocessing pass (Ch.16). Sort by tin, close under adjacent-pair lca, dedup, then a monotone stack links each node to its nearest kept ancestor. Edge weight = depth difference. ASCII only}, label={lst:c30-build}]
// g[] = real tree adjacency; tin[], depth[] known; lca(u,v) from binary lifting
vector<int> buildVT(vector<int> marked){
    sort(marked.begin(), marked.end(), [](int a,int b){ return tin[a]<tin[b]; });
    vector<int> nodes = marked;
    for (int i=0;i+1<(int)marked.size();i++) nodes.push_back(lca(marked[i],marked[i+1]));
    sort(nodes.begin(), nodes.end(), [](int a,int b){ return tin[a]<tin[b]; });
    nodes.erase(unique(nodes.begin(),nodes.end()), nodes.end());   // C = A(S), <= 2k-1

    vector<int> st; st.push_back(nodes[0]);          // stack = current kept root-path
    for (int i=1;i<(int)nodes.size();i++){
        int v = nodes[i];
        while (!isAncestor(st.back(), v)) st.pop_back();   // climb to nearest kept ancestor
        int u = st.back();                              // virtual edge u -> v
        addVEdge(u, v, depth[v]-depth[u]);              // weight = real path length
        st.push_back(v);
    }
    return nodes;   // root of the VT is nodes[0] (the global lca of all marked)
}
\end{lstlisting}

\begin{tracebox}[The stack sweep on \(S=\{4,8,6,10\}\) of the running tree \(T\)]
Sorted by \(\tin\): \(4(3),8(5),6(8),10(10)\).  Adjacent LCAs: \(\lca(4,8)=2\),
\(\lca(8,6)=1\), \(\lca(6,10)=3\).  Closed set sorted by \(\tin\):
\(C=1(1),2(2),4(3),8(5),3(7),6(8),10(10)\) -- \(7=2k-1\) nodes.  Sweep:\\
\(\;\)stack\(=[1]\).  \(v=2\): \(1\) is ancestor of \(2\) \(\Rightarrow\) edge \(1\!\xrightarrow{\;1\;}\!2\), stack\([1,2]\).\\
\(\;\)\(v=4\): \(2\) anc \(4\) \(\Rightarrow\) edge \(2\!\xrightarrow{\;1\;}\!4\), stack\([1,2,4]\).\\
\(\;\)\(v=8\): \(4\) not anc \(8\), pop; \(2\) anc \(8\) \(\Rightarrow\) edge \(2\!\xrightarrow{\;2\;}\!8\) (path \(2\text{-}5\text{-}8\), length \(2\)), stack\([1,2,8]\).\\
\(\;\)\(v=3\): \(8\) no, pop; \(2\) no, pop; \(1\) anc \(3\) \(\Rightarrow\) edge \(1\!\xrightarrow{\;1\;}\!3\), stack\([1,3]\).\\
\(\;\)\(v=6\): \(3\) anc \(6\) \(\Rightarrow\) edge \(3\!\xrightarrow{\;1\;}\!6\), stack\([1,3,6]\).\\
\(\;\)\(v=10\): \(6\) no, pop; \(3\) anc \(10\) \(\Rightarrow\) edge \(3\!\xrightarrow{\;2\;}\!10\) (\(3\text{-}7\text{-}10\)), stack\([1,3,10]\).\\
Six virtual edges, matching the right panel of \(\figref{c30-tree}\); check
\(\dist(8,10)=\depth(8)+\depth(10)-2\depth(\lca(8,10))=3+3-2\cdot0=6\), and the
virtual path \(8\!\to\!2\!\to\!1\!\to\!3\!\to\!10\) has weight \(2+1+1+2=6\). \(\Box\)
\end{tracebox}

\begin{complexitybox}[Construction]
Sorting the \(k\) marked nodes dominates at \(\bigO{k\log k}\); the \(k-1\) adjacent
\(\lca\) calls cost \(\bigO{k\log n}\) with binary lifting (\(\bigO{1}\) with an RMQ
LCA, Ch.~16).  The stack sweep is \(\bigO{k}\) amortised (each node is pushed and
popped once).  Total \(\bigO{k\log k}\) after an \(\bigO{n\log n}\) one-time LCA
preprocess, independent of \(n\) per query -- the whole point. \textbf{Space}
\(\bigO{k}\).
\end{complexitybox}

% ---------------------------------------------------------------
\subsection{Problem Lab: CF 2127E Ancient Tree -- one virtual tree per colour}
% ---------------------------------------------------------------

\solvedproblem{c30-ancient}{Ancient Tree}{CF 2127E}
  {https://codeforces.com/problemset/problem/2127/E}{CF 2100 L5}
  {per-colour virtual tree (sort-by-tin + adjacent-pair LCA); a vertex forced cutie
   iff it is a branch node of 2+ colour virtual trees; DFS colouring to realise the bound}

\begin{problemstatement}[Ancient Tree (statement, constraints and sample verified against the live CF page)]
A rooted tree on \(n\) vertices (root \(1\)); vertex \(i\) has weight \(w_i\) and a
colour \(c_i\in[0,k]\), where \(c_i=0\) means the colour was lost and must be
re-painted with some value in \([1,k]\).  Call vertex \(v\) \textbf{cutie} if there
are two vertices \(x,y\) with \(\lca(x,y)=v\), \(c_x=c_y\), and \(c_x\ne c_v\).  The
cost of the tree is the sum of weights of all cutie vertices.  Repaint every
\(c_i=0\) vertex to minimise the cost, and output the minimum cost together with one
achieving colouring.  \(3\le n\le 2\cdot10^5\), \(2\le k\le n\), \(\sum n\le 2\cdot10^5\)
over tests, \(w_i\le 10^9\).

\textbf{Example (official, 4 tests).} Test 1: \(n=4,k=4\); \(w=5\,5\,5\,5\);
\(c=1\,0\,2\,3\); edges \(1\!-\!2,1\!-\!3,1\!-\!4\).  Output cost \(0\) with a
colouring such as \(1\,4\,2\,3\).  (Full sample set is on the CF page; test 1 is the
warm-up we trace below.)
\end{problemstatement}

\begin{attemptgate}[Attempt it yourself first -- \lv{4}, give yourself 35 minutes]
Focus on the \emph{already-coloured} vertices.  For a fixed colour \(d\), which
vertices are \emph{guaranteed} to become cutie no matter how you paint the blanks?
\peek{a vertex \(v\) is cutie if some colour \(d\ne c_v\) appears at two vertices
whose LCA is \(v\).  Look at the set \(S_d\) of vertices already coloured \(d\).}
\nudge{\(v\) sees a same-colour-\(d\) pair straddling it exactly when \(v\) is an
\(\lca\)-branch of \(S_d\) -- i.e.\ when \(v\) is an internal (adjacent-pair LCA)
node of the \emph{virtual tree of colour \(d\)}.}
\fullhint{Build, for each colour \(d\), the virtual tree of \(S_d\) (sort \(S_d\) by
\(\tin\), add adjacent-pair LCAs).  A vertex that is a branch (non-leaf, LCA) node of
\emph{two or more} colours' virtual trees is unavoidably cutie: whatever colour you
give it, at least one of those colours is different and still straddles it.  A vertex
that is a branch of colour \(d\) and already carries a fixed colour \(c\ne d\) is also
forced.  The answer is the sum of weights over these forced vertices, and a single DFS
that copies the nearest coloured ancestor downwards (painting blanks) achieves
exactly that bound -- no \emph{new} cutie appears because painting a blank with an
ancestor's colour never creates a fresh same-colour straddle.}
\end{attemptgate}

\solstep{Understanding}
The object is not one virtual tree but \(k\) of them -- one per colour class --
and the cost is read off how those trees \emph{overlap}.  A vertex becomes a cutie
precisely when it is an LCA-branch point for a colour different from its own; since
the set of branch points of colour \(d\) is exactly the internal node set of the
virtual tree \(A(S_d)\), the question ``which vertices are forced cutie?'' is the
question ``which vertices appear as branch nodes in \(\ge 2\) colour virtual trees,
or in one whose colour differs from theirs?''  Blanks can then be painted to hit the
bound without adding cost, so the hard part is only \emph{detection} -- a pure
sort-by-tin-plus-adjacent-LCA exercise.

\solstep{Key Observations}
\begin{observation}
A vertex \(v\) is a branch of colour \(d\)'s virtual tree iff, in \(S_d\) sorted by
\(\tin\), \(v=\lca(x_i,x_{i+1})\) for some adjacent pair -- so one pass of adjacent
LCAs over each colour list yields every candidate for free (\(\thmref{c30-size}\)).
\end{observation}
\begin{observation}
If \(v\) is a branch of \(\ge2\) distinct colours, every choice of \(c_v\in[1,k]\)
leaves some colour \(d\ne c_v\) branching at \(v\), so \(v\) is cutie in \emph{all}
valid colourings -- its weight always counts.  If \(v\) is a branch of exactly one
colour \(d\) and \(v\) is blank, paint \(v\) colour \(d\) to dodge the cutie.
\end{observation}

\solstep{The Approach Ladder}
\begin{center}\renewcommand{\arraystretch}{1.25}\setlength{\tabcolsep}{3pt}
\begin{tabular}{@{}c p{0.28\linewidth} p{0.30\linewidth} l c@{}}
\toprule
\textbf{\#} & \textbf{Approach} & \textbf{Idea} & \textbf{Time} & \textbf{Verdict}\\
\midrule
1 & \lv{2}\ try all \(k^{\#\text{blanks}}\) colourings & brute force, check cost each & exponential & \tle \\
2 & \lv{3}\ per-vertex DFS over the whole tree & for each \(v\), scan its subtree for a straddle & \(\bigO{n^2}\) & \tle \\
3 & \lv{4}\ \textbf{per-colour virtual trees} & branch nodes of \(A(S_d)\); count overlaps & \(\bigO{n\log n}\) & \ac \\
\bottomrule
\end{tabular}\end{center}
Rung~2 re-walks \(O(n)\) work at every vertex; rung~3 touches each vertex once per
colour it can possibly branch, and \(\sum_d |S_d|=n\), so the total virtual-tree work
is \(O(n\log n)\).

\solstep{Solution}
Bucket vertices by their fixed colour.  For each colour \(d\) with \(|S_d|\ge2\), sort
\(S_d\) by \(\tin\) and for each adjacent pair compute \(v=\lca(\cdot,\cdot)\): if \(v\)
is blank, record \(d\) as one colour branching at \(v\); if \(v\) is already a colour
\(\ne d\), mark \(v\) forced.  After all colours, a blank vertex is forced cutie iff
\(\ge2\) distinct colours branch there.  Sum weights of forced vertices, then run a
DFS painting each blank with the nearest coloured ancestor (root's colour if none) --
which realises the bound.

\solstep{Implementation}
\begin{lstlisting}[caption={CF 2127E via per-colour virtual trees. bucket by colour, sort each bucket by tin, take adjacent LCAs, and tally the colours branching at each blank vertex; forced = fixed-colour straddle OR >=2 branching colours. ASCII only}, label={lst:c30-ancient}]
#include <bits/stdc++.h>
using namespace std;
const int MAXN=200005;
int n,k,tin[MAXN],depth_,timer,par[19][MAXN];
long long w[MAXN]; int col[MAXN];
vector<int> g[MAXN], byCol[MAXN];
void dfs(int v,int p,int d){ tin[v]=timer++; depth_[v]=d; par[0][v]=p;
  for(int i=1;i<19;i++) par[i][v]=par[i-1][par[i-1][v]];
  for(int u:g[v]) if(u!=p) dfs(u,v,d+1); }
int lca(int u,int v){ if(depth_[u]<depth_[v]) swap(u,v);
  int d=depth_[u]-depth_[v]; for(int i=0;i<19;i++) if(d>>i&1) u=par[i][u];
  if(u==v) return u; for(int i=18;i>=0;i--) if(par[i][u]!=par[i][v]){u=par[i][u];v=par[i][v];}
  return par[0][u]; }
int cntBranch[MAXN]; bool forced[MAXN];
int main(){ int T; scanf("%d",&T); while(T--){
  scanf("%d%d",&n,&k);
  for(int i=1;i<=n;i++){ scanf("%lld",&w[i]); byCol[i].clear(); forced[i]=0; cntBranch[i]=0; }
  for(int i=1;i<=n;i++){ scanf("%d",&col[i]); if(col[i]) byCol[col[i]].push_back(i); }
  for(int i=1;i<n;i++){ int u,v; scanf("%d%d",&u,&v); g[u].push_back(v); g[v].push_back(u); }
  timer=0; dfs(1,0,0);
  for(int d=1;d<=k;d++){ auto &V=byCol[d]; if(V.size()<2) continue;
    sort(V.begin(),V.end(),[](int a,int b){return tin[a]<tin[b];});
    for(int i=0;i+1<(int)V.size();i++){ int x=lca(V[i],V[i+1]);
      if(col[x]) { if(col[x]!=d) forced[x]=1; }   // fixed colour straddled by colour d
      else cntBranch[x]++; } }                     // colour d branches at blank x
  long long ans=0;
  for(int i=1;i<=n;i++){ if(forced[i] || cntBranch[i]>=2) ans+=w[i]; }
  printf("%lld\n",ans);
  // paint blanks with nearest coloured ancestor, then print col[]
  ...
} }
\end{lstlisting}

\solstep{Trace}
\begin{tracebox}[Test 1: a 4-node star, no forced cutie]
Tree: root \(1\) with children \(2,3,4\); colours \(c=1,0,2,3\), all weights \(5\).
Fixed-colour buckets: colour~1\(=\{1\}\), colour~2\(=\{3\}\), colour~3\(=\{4\}\) -- all
size \(1\), so \emph{no} colour has \(\ge2\) vertices and the adjacent-pair loop never
runs.  No vertex is a branch of any colour, so \(\mathit{forced}=\emptyset\),
\(\mathit{cntBranch}=0\) everywhere, answer \(\boxed{0}\).  Paint the single blank
\(c_2\) with any colour \(<\!\!\to\) the sample uses \(4\); a colour appearing at just
one vertex cannot straddle anything, so cost stays \(0\).  Matches the official output.
\end{tracebox}

\begin{warningbox}[Three ways the detection goes wrong]
\begin{enumerate}[nosep,leftmargin=1.6em]
  \item Forgetting the \emph{fixed}-colour straddle: a branch node \(x\) with a real
        colour \(c_x\ne d\) is forced even though only colour \(d\) branches there --
        you cannot repaint \(x\).  The \(\mathit{forced}[x]\) flag, not just
        \(\mathit{cntBranch}\ge2\), is what catches it.
  \item Sorting the bucket by vertex id instead of \(\tin\) breaks the
        adjacent-pair-LCA lemma (\(\thmref{c30-size}\)) and misses branch nodes.
  \item Painting blanks top-down must copy the \emph{nearest coloured ancestor};
        painting with a colour that also appears deeper in a different subtree can
        manufacture a cutie the bound said was avoidable.
\end{enumerate}
\end{warningbox}

\begin{hardeningbox}
\begin{itemize}[nosep,leftmargin=1.4em]
  \item \textbf{Only build a virtual tree when \(|S_d|\ge2\).} A one-vertex colour has
        no adjacent pair and cannot branch; skipping it is both correct and faster.
  \item \textbf{\(\sum_d|S_d|=n\).} Total adjacent-LCA work over all colours is
        \(\bigO{n\log n}\); the per-colour sort is the cost, not any single colour.
  \item \textbf{Weights are \(10^9\), sums overflow \texttt{int}.} Accumulate the cost
        in \texttt{long long}.
  \item \textbf{Ties are free.} The statement accepts \emph{any} minimum-cost
        colouring, so paint greedily once the cost is known.
\end{itemize}
\end{hardeningbox}

\begin{recallbox}[26.1 in one breath]
\emph{Close the marked set by sorting it by \(\tin\) and adding the \(k-1\) adjacent
\(\lca\)s; dedup gives at most \(2k-1\) nodes (\(\thmref{c30-size}\)).  Link them with a
monotone stack whose predicate is Ch.~5's ancestor test, and give each edge weight
\(\depth(v)-\depth(u)\) (\(\coderef{c30-build}\)).  CF 2127E is that construction
run once per colour: a vertex is forced cutie when it is a branch node of two colours'
virtual trees (or of one whose fixed colour differs), so the answer is a sum over
\emph{overlapping} virtual trees -- \(\bigO{n\log n}\), never a per-vertex subtree scan.}
\selftest{
  \item Why do \(k-1\) adjacent-pair LCAs close the set, and not \(\binom{k}{2}\)? \emph{(A branch splits marked nodes into two consecutive \(\tin\)-buckets; the split is visible between neighbours. \(\thmref{c30-size}\).)}
  \item What is the size ceiling of \(A(S)\)? \emph{\(2k-1\).}
  \item In 2127E, when is a blank vertex forced cutie? \emph{When it is a branch node of \(\ge2\) colour virtual trees (Observation~2).)}
}
\end{recallbox}
